\cleardoublepage\pagestyle{empty}\bfseries
\begin{center}
    \zihao{3} 毕业论文（设计）文献综述和开题报告考核
\end{center}
% 将该标题加入目录（在章节中插入）
% \addcontentsline{toc}{section}{\hspace{-2.2em} 毕业论文（设计）文献综述和开题报告考核}
\ \\
\zihao{4}对文献综述、外文翻译和开题报告评语及成绩评定
\vfill
该三合一材料整体完成质量高，显示出张三同学已进行了深入、系统的文献调研，对本毕业设计课题的研究背景、科学意义、技术路线及实施方案有了清晰且准确的理解。文献综述部分内容详实，脉络清晰；开题报告部分目标明确，方案具体，规划合理；外文翻译译文准确、通顺，专业术语翻译得当。材料结构完整，逻辑清晰，表明该同学已具备开展本项研究所需的基本知识储备和初步的科研规划能力。

希望接下来张三同学能严格按照技术路线和研究计划开展实验工作，在具体实践中深化对理论的理解，锻炼解决复杂工程科学问题的能力，并注意及时、规范地记录实验过程和数据处理结果，为最终的毕业论文撰写奠定坚实基础。

\begin{flushright}
\begin{table}[!htbp]
    \renewcommand{\arraystretch}{3.5}  % 调整行高，可以增大或减小数值
    \hfill
    \begin{tabular}{|>{\centering\arraybackslash}m{3cm}|>{\centering\arraybackslash}m{2.5cm}|>{\centering\arraybackslash}m{2.5cm}|>{\centering\arraybackslash}m{2.5cm}|} \hline
        \bfseries\CJKfamily{Songti}\zihao{4}成绩比例 & \bfseries\CJKfamily{Songti}\zihao{5}文献综述占（10\%） & \bfseries\CJKfamily{Songti}\zihao{5}开题报告占（15\%）& \bfseries\CJKfamily{Songti}\zihao{5}外文翻译占（5\%）\\ \hline
        \bfseries\CJKfamily{Songti}\zihao{4}分值 &  &  & \\ \hline
    \end{tabular}
\end{table}  
{\hfill {\zihao{-4}\bfseries 开题报告答辩小组负责人（签名）\underline{\makebox[3cm]{}}}}\\
\hfill {\zihao{-4}  2026年\quad\quad  月\quad\quad  日}
\end{flushright} 